%%%%%%%%%%%%%%%%%%%%%%%%%%%%%%%%%%%%%%%%%
% Beamer Presentation
% LaTeX Template
% Version 1.0 (10/11/12)
%
% This template has been downloaded from:
% http://www.LaTeXTemplates.com
%
% License:
% CC BY-NC-SA 3.0 (http://creativecommons.org/licenses/by-nc-sa/3.0/)
%
% Modified by Jeremie Gillet in November 2015 to make an OIST Skill Pill template
%
%%%%%%%%%%%%%%%%%%%%%%%%%%%%%%%%%%%%%%%%%

%----------------------------------------------------------------------------------------
%	PACKAGES AND THEMES
%----------------------------------------------------------------------------------------

\documentclass{beamer}

\mode<presentation> {

\usetheme{Madrid}

\definecolor{OISTcolor}{rgb}{0.65,0.16,0.16}
\usecolortheme[named=OISTcolor]{structure}

%\setbeamertemplate{footline} % To remove the footer line in all slides uncomment this line
%\setbeamertemplate{footline}[page number] % To replace the footer line in all slides with a simple slide count uncomment this line

\setbeamertemplate{navigation symbols}{} % To remove the navigation symbols from the bottom of all slides uncomment this line
}
\setbeamertemplate{itemize items}[default]
\setbeamertemplate{enumerate items}[default]
\usepackage{graphicx} % Allows including images
\usepackage{booktabs} % Allows the use of \toprule, \midrule and \bottomrule in tables
\usepackage{textpos} % Use for positioning the Skill Pill logo
\usepackage{fancyvrb}
\usepackage{tikz}
\usepackage{hyperref}
\usepackage{listings}
\usepackage{color}

\definecolor{dkgreen}{rgb}{0,0.6,0}
\definecolor{gray}{rgb}{0.5,0.5,0.5}
\definecolor{mauve}{rgb}{0.58,0,0.82}

\lstset{frame=tb,
  language=python,
  aboveskip=3mm,
  belowskip=3mm,
  showstringspaces=false,
  columns=flexible,
  basicstyle={\small\ttfamily},
  numbers=none,
  numberstyle=\tiny\color{gray},
  keywordstyle=\color{blue},
  commentstyle=\color{dkgreen},
  stringstyle=\color{mauve},
  breaklines=true,
  breakatwhitespace=true,
  tabsize=3
}

\newcommand*{\lstitem}[1]{
  \setbox0\hbox{\lstinline{#1}}  
  \item[\usebox0]  
  % \item[\hbox{\lstinline{#1}}]
  \hfill \\
}

%----------------------------------------------------------------------------------------
%	TITLE PAGE
%----------------------------------------------------------------------------------------

\title[Skill Pill]{Skill Pill: Julia} % The short title appears at the bottom of every slide, the full title is only on the title page
\subtitle{Lecture 4: Distributed and parallel computing}

\author{Valentin Churavy \and James Schloss} % Your name
\institute[OIST] % Your institution as it will appear on the bottom of every slide, may be shorthand to save space
{
Okinawa Institute of Science and Technology \\ % Your institution for the title page
\textit{valentin.churavy@oist.jp} \\
\textit{james.schloss@oist.jp} % Your email address
}
\date{July 13, 2017} % Date, can be changed to a custom date

\begin{document}

\setbeamertemplate{background}{\includegraphics[width=\paperwidth]{SPbackground.png}} % Adding the background logo

\begin{frame}
\vspace*{1.4cm}
\titlepage % Print the title page as the first slide
\end{frame}

\setbeamertemplate{background}{} % No background logo after title frame

\addtobeamertemplate{frametitle}{}{% Adding the Skill Pill logo on the title screen after title frame
\begin{textblock*}{100mm}(.92\textwidth,-0.9cm)
\includegraphics[height=0.85cm]{julia.pdf}
\end{textblock*}}

\begin{frame}
  \tableofcontents
\end{frame}

\section{Levels of parallelism}
\begin{frame}{Introduction}
  \pause
  \begin{block}{Necessary packages}
    \begin{itemize}
      \item SIMD.jl
      \item MPI.jl
      \item DistributedArrays.jl
      \item CUDAnative.jl if your computer has a NVidia GPU
    \end{itemize}
  \end{block}
\end{frame}
\begin{frame}{Levels of parallelism}
  \begin{itemize}
    \item Instruction level parallelism
    \item Shared-memory and threading
    \item Distributed
    \item Accelerators e.g.GPGPU
  \end{itemize}
\end{frame}

\section{Instruction level parallelism}
\begin{frame}[fragile]{What is instruction level parallelism}

  \begin{lstlisting}
  function padd(a, b, x, y)
    c = a + b
    z = x + y
    return c, z
  end
  \end{lstlisting}

  \begin{block}{Observation}
     The two operations are independent of each other and we could execute them in parallel.
  \end{block}

  \begin{enumerate}
    \item Use \lstinline{@code_llvm} and \lstinline{@code_native} to understand what is happening
    \item Establish a baseline performance with \lstinline{@benchmark}
    \item Start Julia with \lstinline{julia -O3}
    \item Compare the llvm and native code and your benchmark results
    \item Note that there is next to no performance benefit in this example, but that changes once you scale up
  \end{enumerate}
\end{frame}

\begin{frame}[fragile]{SIMD and loops}
  \begin{lstlisting}
  function add!(out, a, b)
    @assert length(a) == length(b)
    @assert length(a) == length(out)

    for i in 1:length(a)
       out[i] = a[i] + b[i]
    end
  end
  \end{lstlisting}
  \begin{block}{Observation}
    Each loop iteration is independent from each other.
  \end{block}
  \begin{enumerate}
    \item Learn about \lstinline{@inline} and \lstinline{@simd}
    \item Note that LLVM will vectorise this loop even without \lstinline{@simd}
  \end{enumerate}
\end{frame}
\begin{frame}[fragile]{Explicit SIMD}
  \begin{block}{SIMD.jl}
    Instead of relying on the compiler to optimise and vectorise our code correctly we can also write explicit SIMD code.
  \end{block}
  \begin{lstlisting}
  using SIMD
  function add(out::Vector{Float64}, x::Vector{Float64}, y::Vector{Float64})
    # My laptop supports AVX 256bit 4xFloat64
    @assert length(x) % 4 == 0
    for i in 1:4:length(x)
      vx = vload(Vec{4, Float64}, x, i)
      vy = vload(Vec{4, Float64}, y, i)
      vo = vx + vy
      vstore(vo, out, i)
    end
  end
  \end{lstlisting}
\end{frame}
\section{Threading}
\begin{frame}
  Explain current fork join model and caveats
\end{frame}
\begin{frame}
  Simple example
\end{frame}
\begin{frame}
  complex example with loop splitting, random etc...
\end{frame}
\begin{frame}
  Atomics
\end{frame}
\section{Distributed}
\begin{frame}
  MPI
\end{frame}
\begin{frame}
  DistibutedArrays
\end{frame}
\section{CUDAnative}
\end{document}
